\documentclass[submission,copyright,creativecommons]{eptcs}
\providecommand{\event}{PCC 2026: Proof, Computation, and Complexity}
\usepackage{underscore}

\title{O(1) Memory Validation:\\ Physical Memory Caching for Deterministic Execution Verification}
\author{Rafael D. De Paz
\institute{Vanguard Research Array \\ Sovereign Architecture Group \\ Las Vegas, NV}
\email{rafael@rdepaz.com}
}
\def\titlerunning{O(1) Physical Validation Schema}
\def\authorrunning{R. D. De Paz}


\usepackage[top=1in, bottom=1in, left=1in, right=1in]{geometry}
\usepackage{draftwatermark}
\SetWatermarkText{SOVEREIGN PROTOCOL // ID: AB438ED}
\SetWatermarkScale{0.5}
\SetWatermarkLightness{0.94}
\begin{document}
\maketitle

\begin{abstract}
As computation scales across probabilistic and distributed ledgers, verifying state transitions requires linear $O(N)$ checksum computation, creating a critical bottleneck in high-frequency validation nodes. We introduce a novel mathematical architecture that decouples structural integrity checking from payload parsing by utilizing physical thermodynamic memory caching techniques. By anchoring dynamic state changes locally into immutable, physical filesystem tethers linked via continuous SHA-256 state arrays, we formalize a framework capable of verifying complex deterministic execution environments with a constant time $O(1)$ complexity constraint.
\end{abstract}

\end{document}